\documentclass[conference]{IEEEtran}
\IEEEoverridecommandlockouts

\usepackage[utf8]{inputenc}
\usepackage[T1]{fontenc}
\usepackage[brazil]{babel}
\usepackage{cite}
\usepackage{amsmath,amssymb,amsfonts}
\usepackage{graphicx}
\usepackage{textcomp}
\usepackage{xcolor}
\usepackage{booktabs}
\usepackage{multirow}
\usepackage{url}
\usepackage{hyperref}
\usepackage{tikz}
\usetikzlibrary{shapes.geometric, arrows.meta, positioning, fit, backgrounds}

\hypersetup{
    colorlinks=true,
    linkcolor=black,
    citecolor=black,
    urlcolor=blue
}

\begin{document}

\title{ Validação Discriminativa Calibrada na Conversão de Requisitos em Linguagem Natural para Templates MASTeR: Uma Abordagem em Cascata com JEV }

\author{
\IEEEauthorblockN{Lucas F. B. M.}
\IEEEauthorblockA{\textit{Disciplina de Análise Avançada de Requisitos} \\
\textit{Departamento de Computação}\\
Brasil \\
lucasfbm@exemplo.edu.br}
}

\maketitle

\begin{abstract}
\textbf{[Contexto e Motivação]} Templates sintáticos semiformais, como o sistema MASTeR, reduzem ambiguidades e impõem atomicidade a especificações de requisitos em linguagem natural (LN). Recentemente, Ro\ss kothen et al. (IEEE RE 2026) demonstraram que Grandes Modelos de Linguagem (LLMs) inseridos em um ciclo reflexivo (\textit{Reflection Loop}) alcançam qualidade próxima à humana na conversão em lote de requisitos livres para o padrão MASTeR. \textbf{[Questão/Problema]} Contudo, o ciclo reflexivo proposto na literatura apoia-se em validadores generativos (\textit{LLM-as-a-Judge}) e filtros auxiliares que exigem até $2k + 1$ chamadas autorregressivas por iteração para $k$ requisitos candidatos. Além de elevar o custo financeiro e a latência em projetos industriais, a ausência de probabilidades calibradas provoca edições desnecessárias (\textit{over-editing}) em requisitos simples e falha em sinalizar requisitos complexos que demandam supervisão humana. \textbf{[Ideias Principais e Resultados]} Propomos a arquitetura \textit{JEV-Cascade Reflection Loop}, que substitui os validadores generativos e o filtro de acionabilidade por uma única inferência paralela \textit{System One} utilizando o modelo discriminativo JEV. A abordagem avalia simultaneamente escolhas tipadas de sub-template, \textit{checklists} booleanos condicionados ao requisito fonte e rubricas ordinais com confiança calibrada, roteando cada caso entre aceitação rápida (\textit{Fast Accept}), reescrita determinística (\textit{Guided Rewrite}) ou escalonamento humano (\textit{Escalate}). \textbf{[Contribuição]} Apresentamos a formulação arquitetural do \textit{JEV-Cascade}, a eliminação por construção do estágio de filtragem de feedback por LLM e um pacote experimental reprodutível avaliado sobre o conjunto público de 249 requisitos fonte (410 requisitos MASTeR de referência), contrastando eficácia sintático-semântica, latência, custo por lote e calibração de escalonamento em múltiplas famílias de LLMs geradores.
\end{abstract}

\begin{IEEEkeywords}
Engenharia de Requisitos, Templates de Requisitos, MASTeR, Grandes Modelos de Linguagem, LLM-as-a-Judge, Modelos de Decisão System One, JEV.
\end{IEEEkeywords}

\section{Introdução}

Requisitos escritos em linguagem natural (LN) irrestrita permanecem o veículo predominante para comunicação entre partes interessadas em projetos industriais de software e sistemas embarcados~\cite{darif2026controlled}. Entretanto, a flexibilidade inerente à prosa livre frequentemente introduz ambiguidades lexicais, omissão de atores, voz passiva e agrupamento de múltiplas obrigações independentes em uma única sentença, prejudicando a rastreabilidade, a verificação formal e a garantia de qualidade em escala~\cite{incose2023guide, grosser2024benchmarking}. Para mitigar esses problemas sem exigir linguagens lógicas matemáticas inacessíveis a clientes não técnicos, sistemas de templates semiformais --- como EARS e MASTeR (\textit{Muster für Anforderungen und System-Tests im Engineering} / SOPHIST)~\cite{sophist2020master, rupp2014anforderungsschablonen} --- estabelecem moldes sintáticos fixos que elevam a clareza, a singularidade e a conformidade normativa~\cite{grosser2023comparative, grosser2024benchmarking}.

Apesar de seus benefícios comprovados em domínios críticos de segurança (como os setores automotivo e aeroespacial), migrar especificações legadas contendo milhares de requisitos em texto livre para o padrão MASTeR impõe alto custo manual: exige domínio simultâneo do negócio e da gramática de sub-templates, além da decomposição (\textit{splitting}) frequente de requisitos compostos em múltiplos requisitos atômicos~\cite{rosskothen2026converting}. Recentemente, Ro\ss kothen et al.~\cite{rosskothen2026converting} apresentaram uma arquitetura orientada à indústria para automatizar essa conversão em lote utilizando Grandes Modelos de Linguagem (LLMs) organizados em um ciclo reflexivo (\textit{Reflection Loop}~\cite{ji2023mitigating}), combinando um validador sintático MASTeR, um avaliador alinhado ao guia INCOSE~\cite{incose2023guide} e um filtro de acionabilidade de feedback.

Uma análise detalhada da arquitetura e do pacote de replicação público de Ro\ss kothen et al.~\cite{rosskothen2026zenodo}, contudo, expõe três limitações críticas para adoção em larga escala:
\begin{enumerate}
    \item \textbf{Sobrecarga multiplicativa de chamadas generativas:} Para cada requisito fonte decomposto em $k$ requisitos candidatos, uma única iteração de validação executa até $2k + 1$ chamadas autorregressivas a LLMs juízes (\textit{Template Validator}, \textit{INCOSE Validator} em nove categorias e \textit{Feedback Filter}), reenviando repetidamente prompts extensos (~1.400 tokens) apenas para produzir justificativas textuais intermediárias;
    \item \textbf{Paradoxo de \textit{over-editing} sem calibração:} Nos dados empíricos publicados pelos próprios autores~\cite{rosskothen2026zenodo}, o modelo \texttt{Qwen3-235B} sem validadores generativos (\texttt{minimal-qwen235}) obteve 68\% de avaliações humanas positivas (categorias C1/C2) e pontuações automáticas ligeiramente superiores às da configuração completa com juízes (\texttt{full-qwen235}: \texttt{CorpBERTF1} de $0{,}834$ contra $0{,}823$; \textit{Compound Score} de $0{,}811$ contra $0{,}797$), indicando que ciclos reflexivos não calibrados podem degradar rascunhos inicialmente corretos;
    \item \textbf{Falta de roteamento explícito por incerteza:} Embora os autores demonstrem que tanto os LLMs quanto as métricas automáticas sofrem forte queda de desempenho em requisitos complexos (\textit{split factor} $\ge 2$ ou mais de 30 palavras), o fluxo original não fornece probabilidades calibradas para decidir automaticamente quando aceitar um requisito, quando reescrevê-lo ou quando escalá-lo para inspeção humana.
\end{enumerate}

Neste artigo, investigamos a substituição da pilha de validadores generativos (\textit{LLM-as-a-Judge}) por um modelo discriminativo de decisão rápida (\textit{System One}), especificamente o \textbf{JEV}~\cite{typesafe2026jev}. Diferentemente de LLMs autorregressivos que decodificam tokens de texto sequencialmente para justificar notas, o JEV avalia um estado estruturado contra múltiplas consultas tipadas (\texttt{Choice}, \texttt{Noul} e \texttt{Score}) em um único passe paralelo, retornando distribuições de probabilidade calibradas em milissegundos. Combinando validação discriminativa no ciclo reflexivo com o padrão arquitetural \textit{Accept-Rewrite-Escalate}, propomos o \textbf{JEV-Cascade Reflection Loop}.

As principais contribuições deste trabalho são:
\begin{itemize}
    \item \textbf{(i)} O projeto arquitetural do \textit{JEV-Cascade Reflection Loop}, que condiciona a validação simultaneamente ao requisito fonte $S$ e ao candidato $R_i$, eliminando por construção o estágio extra de filtragem de feedback por LLM e mapeando falhas para instruções determinísticas de reparo;
    \item \textbf{(ii)} Um mecanismo de roteamento triplo baseado em probabilidades calibradas que protege rascunhos de alta confiança contra \textit{over-editing} (\textit{Fast Accept}), aciona reescrita apenas diante de violações claras (\textit{Guided Rewrite}) e isola casos de incerteza epistêmica para engenheiros de requisitos (\textit{Escalate});
    \item \textbf{(iii)} Um protocolo e repositório experimental reprodutível avaliado sobre os 249 requisitos fonte e 410 requisitos MASTeR de referência~\cite{grosser2023comparative, rosskothen2026zenodo}, contrastando eficácia sintático-semântica, latência e custo operacional em modelos abertos (\texttt{Qwen3-235B}) e comerciais (\texttt{Gemini 3.5 Flash}).
\end{itemize}

\section{Fundamentação Teórica e Trabalhos Relacionados}

\subsection{Sistema de Templates MASTeR e Diretrizes INCOSE}
O sistema MASTeR~\cite{sophist2020master, rupp2014anforderungsschablonen} estrutura requisitos em linguagem controlada combinando um \textbf{template principal} obrigatório e um \textbf{sub-template de condição} opcional, vinculados a verbos modais estritos (\textit{shall}, \textit{should} ou \textit{will}). Os templates principais dividem-se em requisitos funcionais (\textit{FunctionMASTeR}, subdividido em atividade independente do sistema, interação com usuário e requisito de interface) e não funcionais (\textit{PropertyMASTeR}, \textit{EnvironmentMASTeR} e \textit{ProcessMASTeR}). Quando aplicável, a cláusula condicional precede o template principal segundo três padrões: \textit{LogicMASTeR} (iniciado por \textit{``If \dots''}), \textit{EventMASTeR} (\textit{``As soon as \dots''}) ou \textit{TimeMASTeR} (\textit{``As long as \dots''}), exigindo repetição explícita da palavra-chave condicional após operadores lógicos \textit{and}/\textit{or}.

Complementarmente à estrutura sintática, o \textit{INCOSE Guide to Writing Requirements}~\cite{incose2023guide} estabelece nove características individuais para requisitos bem formados (C1--C9): \textit{necessário}, \textit{apropriado}, \textit{não ambíguo}, \textit{completo}, \textit{singular}, \textit{viável}, \textit{verificável}, \textit{correto} e \textit{conforme}. Estudos empíricos de Gro\ss er et al.~\cite{grosser2023comparative, grosser2024benchmarking} confirmam que o uso do MASTeR eleva significativamente a singularidade (atomicidade) e a conformidade normativa frente ao texto livre, embora exija treinamento especializado para redação manual.

\subsection{Conversão Automática de Requisitos com LLMs}
Abordagens clássicas de reescrita automática de requisitos baseavam-se em regras gramaticais manuais, ontologias e árvores sintáticas~\cite{arora2015automated}. Com o advento dos LLMs, Norheim e Rebentisch~\cite{norheim2024structuring} exploraram preenchimento de lacunas (\textit{slot-filling}) via \textit{one-shot prompting} para o template EARS, porém sem avaliação sistemática em larga escala.

Ro\ss kothen et al.~\cite{rosskothen2026converting} avançaram o estado da arte ao propor a geração integral de sentenças MASTeR com decomposição livre ($1 \rightarrow n$ requisitos atômicos) controlada por decodificação estruturada (\textit{Constrained Decoding}) e refinada em um \textit{Reflection Loop}~\cite{ji2023mitigating}. Avaliando quatro configurações sobre 249 requisitos de cinco projetos reais reescritos por humanos em 410 sentenças MASTeR~\cite{grosser2023comparative}, os autores mostraram que métricas baseadas em \textit{BERTScore}~\cite{zhang2020bertscore}, \textit{Embedding Mover's Distance} (EMD, derivada de WMD~\cite{kusner2015word}) e conformidade de sub-templates apresentam correlação positiva moderada ($\rho = 0{,}503$ combinada) com o julgamento de especialistas da indústria.

\subsection{LLM-as-a-Judge Generativo vs. Modelos Discriminativos System One (JEV)}
No paradigma convencional de \textit{LLM-as-a-Judge}, um modelo generativo recebe instruções extensas e produz cadeias de raciocínio ou justificativas textuais seguidas de um veredito booleano ou numérico. Embora flexível, essa estratégia sofre de três gargalos em pipelines de Engenharia de Requisitos em lote: (i) \textbf{alta latência e custo de geração}, pois a decodificação autorregressiva é sequencial por token; (ii) \textbf{variância e deriva de justificativas}, exigindo estágios extras de filtragem para impedir que críticas excessivas induzam alucinações no modelo gerador; e (iii) \textbf{má calibração probabilística}, dificultando a definição de limiares confiáveis de aceitação automática.

Em contrapartida, modelos discriminativos de decisão denominados \textit{System One}, exemplificados pela arquitetura \textbf{JEV} (\textit{TypeSafe AI})~\cite{typesafe2026jev}, separam o julgamento estruturado da geração de prosa. Dado um contexto de entrada (\texttt{state}), o JEV avalia múltiplas consultas tipadas em um único passe paralelo (\textit{forward pass}), sem gerar tokens de texto livre:
\begin{itemize}
    \item \texttt{Choice}: classificação multiclasse mutuamente exclusiva retornando a classe escolhida, a distribuição de probabilidades sobre todas as opções e um índice de confiança calibrado;
    \item \texttt{Noul}: verificação binária de condição (\textit{checklist}), retornando diretamente a probabilidade calibrada $p \in [0, 1]$ de atendimento à regra (onde $p \approx 0{,}5$ indica alta incerteza epistêmica);
    \item \texttt{Score}: avaliação ordinal sobre níveis descritivos ordenados, retornando a posição contínua ponderada pelas probabilidades de cada nível e sua respectiva confiança.
\end{itemize}
Essa distinção arquitetural permite executar auditorias multi-critério em dezenas a poucas centenas de milissegundos, viabilizando políticas formais de cascata (\textit{Accept or Escalate}) em fluxos industriais de requisitos.

\section{Abordagem Proposta: JEV-Cascade Reflection Loop}

\subsection{Análise de Gargalos do Ciclo Reflexivo Base}
Na implementação de referência de Ro\ss kothen et al.~\cite{rosskothen2026converting, rosskothen2026zenodo}, após a geração inicial produzir $k$ requisitos candidatos $R = \{R_1, \dots, R_k\}$ a partir do requisito fonte $S$, cada iteração do ciclo reflexivo executa:
\begin{enumerate}
    \item até $k$ chamadas ao \textit{Template Validator} generativo (acionado em configuração híbrida sempre que o analisador gramatical livre de contexto identifica falha estrutural), retransmitindo a especificação completa do formato MASTeR;
    \item $k$ chamadas ao \textit{INCOSE Validator}, gerando notas e comentários textuais para cada uma das 9 categorias (C1--C9) avaliando $R_i$ isoladamente (sem acesso a $S$);
    \item $1$ chamada ao \textit{Feedback Filter}, que confronta os comentários gerados contra $S$ para descartar exigências impossíveis de cumprir sem alucinação (por exemplo, exigir uma tolerância em milissegundos inexistente no texto original).
\end{enumerate}
Portanto, uma única iteração consome até $2k + 1$ chamadas de validação generativa antes mesmo da chamada de reescrita. Para um requisito com fator de divisão $k = 3$ ao longo de duas iterações, o pipeline acumula até $1 + 2 \times (2(3) + 1 + 1) = 17$ invocações autorregressivas completas.

\begin{figure}[t]
\centering
\resizebox{\columnwidth}{!}{%
\begin{tikzpicture}[
    node distance=0.65cm and 0.6cm,
    box/.style={rectangle, rounded corners=3pt, draw=black!80, thick, fill=blue!5, align=center, minimum height=0.8cm, font=\footnotesize, inner sep=5pt},
    jevbox/.style={rectangle, rounded corners=3pt, draw=teal!80!black, thick, fill=teal!8, align=center, minimum height=0.85cm, font=\footnotesize, inner sep=5pt},
    decision/.style={diamond, aspect=2.1, draw=orange!80!black, thick, fill=orange!10, align=center, font=\scriptsize\bfseries, inner sep=2pt},
    outaccept/.style={rectangle, rounded corners=3pt, draw=green!60!black, thick, fill=green!10, align=center, font=\scriptsize\bfseries, inner sep=4pt},
    outesc/.style={rectangle, rounded corners=3pt, draw=red!70!black, thick, fill=red!8, align=center, font=\scriptsize\bfseries, inner sep=4pt},
    arrow/.style={-{Stealth[length=2.2mm]}, thick, draw=black!75}
]
    \node[box] (src) {Requisito Fonte em LN ($S$)};
    \node[box, below=0.5cm of src] (gen) {\textbf{LLM Gerador} (\textit{Constrained JSON})\\Formalize \& Split: $R = \{R_1, \dots, R_k\}$};
    \node[jevbox, below=0.55cm of gen] (jev) {\textbf{JEV System One} ($1$ passe paralelo por $R_i$ com estado $(S, R_i)$)\\
    \texttt{Choice} (Template/Cond.) + $6\times$\texttt{Noul} Acionáveis + \texttt{Score} (C4..C1)};
    \node[decision, below=0.6cm of jev] (router) {Roteador\\Calibrado};

    \node[outaccept, below left=0.75cm and -0.4cm of router] (accept) {PATH A: Fast Accept\\($\forall p_j \ge \tau_{\text{acc}}$)\\0 chamadas extras};
    \node[outesc, below right=0.75cm and -0.4cm of router] (escalate) {PATH C: Escalate Humano\\($\tau_{\text{rej}} \le p_j < \tau_{\text{acc}}$ ou iter $=M$)\\Fila de Revisão};
    \node[box, fill=yellow!12, left=0.65cm of jev, text width=2.5cm] (repair) {\textbf{PATH B: Mapa Determinístico}\\Chave reprovada ($p_j < \tau_{\text{rej}}$) $\rightarrow$ Instrução de reparo};

    \draw[arrow] (src) -- (gen);
    \draw[arrow] (gen) -- (jev);
    \draw[arrow] (jev) -- (router);
    \draw[arrow] (router) -| node[pos=0.25, above, font=\scriptsize] {Aprovado} (accept);
    \draw[arrow] (router) -| node[pos=0.25, above, font=\scriptsize] {Incerto} (escalate);
    \draw[arrow] (router.west) -| node[pos=0.22, below, font=\scriptsize] {Violação clara} (repair.south);
    \draw[arrow] (repair.north) |- (gen.west);
\end{tikzpicture}%
}
\caption{Visão arquitetural do \textit{JEV-Cascade Reflection Loop} proposto, unificando validação sintático-semântica condicionada à fonte e roteamento triplo calibrado.}
\label{fig:architecture}
\end{figure}

\subsection{Validação System One Condicionada à Fonte}
A Figura~\ref{fig:architecture} ilustra o \textit{JEV-Cascade Reflection Loop}. Mantemos inalterados o prompt de sistema MASTeR e o schema estruturado de geração de Ro\ss kothen et al.~\cite{rosskothen2026converting} para preservar estritamente a capacidade generativa e a comparabilidade experimental. Quando a ação escolhida pelo gerador é diferente de \texttt{no\_update}, cada requisito candidato $R_i$ é submetido ao JEV juntamente com o texto fonte original $S$, compondo o estado $\text{state} = (S, R_i)$.

Em uma única requisição paralela \texttt{system\_one}, o JEV avalia nove primitivas tipadas:
\begin{itemize}
    \item \textbf{Duas primitivas \texttt{Choice}:} classificam o template principal (\texttt{FunctionMASTeR} nas três variantes, \texttt{PropertyMASTeR}, \texttt{EnvironmentMASTeR}, \texttt{ProcessMASTeR} ou \texttt{NonConforming}) e o sub-template condicional (\texttt{LogicMASTeR}, \texttt{EventMASTeR}, \texttt{TimeMASTeR}, \texttt{No Condition} ou \texttt{InvalidCondition}), fornecendo probabilidades calibradas e confiança;
    \item \textbf{Seis primitivas \texttt{Noul} estritamente acionáveis:} (1) \texttt{has\_modal\_verb} (presença correta de \textit{shall/should/will}); (2) \texttt{conforms\_to\_master\_syntax} (aderência à ordem de slots e encadeamento condicional); (3) \texttt{is\_singular\_atomic} (ausência de obrigações múltiplas); (4) \texttt{has\_explicit\_actor\_or\_system} (sujeito explícito sem pronomes vagos); (5) \texttt{preserves\_source\_details} (retenção fiel de qualificadores presentes em $S$); e (6) \texttt{no\_hallucinated\_values} (ausência de números ou unidades inventadas fora de $S$);
    \item \textbf{Uma primitiva \texttt{Score}:} rubrica ordinal de quatro níveis alinhada às categorias práticas C4 a C1 do estudo original ($s \in [0, 3]$).
\end{itemize}

\textbf{Eliminação por construção do Filtro de Feedback:} Como todas as primitivas \texttt{Noul} observam simultaneamente $S$ e $R_i$ e restringem-se a propriedades gramaticais ou informações presentes em $S$, o validador jamais solicita detalhes ausentes no documento fonte. Assim, cada primitiva $j$ reprovada mapeia deterministicamente por tabela hash para uma instrução concisa de correção, eliminando 100\% das chamadas de LLM do \textit{Feedback Filter} e o custo de geração de comentários livres.

\subsection{Política de Roteamento Triplo Calibrado (Cascade)}
Sejam $\tau_{\text{rej}}$ e $\tau_{\text{acc}}$ os limiares inferior e superior de probabilidade ($0 < \tau_{\text{rej}} < \tau_{\text{acc}} < 1$, configurados por padrão em $\tau_{\text{rej}} = 0{,}40$ e $\tau_{\text{acc}} = 0{,}75$). Para o conjunto de probabilidades $\{p_1, \dots, p_6\}$ das verificações \texttt{Noul} e o escore ordinal $s$ de um candidato $R_i$, o roteador define três rotas mutuamente exclusivas:
\begin{enumerate}
    \item \textbf{Path A --- \textit{Fast Accept}:} Se $\min_j p_j \ge \tau_{\text{acc}}$, os templates classificados não são inválidos e $s \ge 1{,}8$, o requisito é aceito imediatamente sem reescrita, protegendo boas gerações contra \textit{over-editing} e zerando o custo iterativo;
    \item \textbf{Path B --- \textit{Deterministic Guided Rewrite}:} Se existe ao menos um critério com falha clara ($p_j < \tau_{\text{rej}}$ ou classe \texttt{NonConforming} com alta confiança) e o número de iterações é inferior a $M$, as instruções determinísticas correspondentes às chaves violadas alimentam o próximo turno do LLM gerador;
    \item \textbf{Path C --- \textit{Uncertainty Escalation}:} Se não há violação clara ($\forall j, p_j \ge \tau_{\text{rej}}$), mas ao menos um critério recai na faixa de incerteza epistêmica ($\tau_{\text{rej}} \le p_j < \tau_{\text{acc}}$) ou as $M$ tentativas de reescrita se esgotam sem atingir \textit{Fast Accept}, o requisito é marcado como \texttt{escalated\_to\_human = True} para revisão prioritária por especialistas.
\end{enumerate}

\section{Desenho Experimental}

Com o objetivo de avaliar rigorosamente o impacto da validação discriminativa calibrada frente ao paradigma generativo tradicional, estruturamos nossa investigação em torno de três Questões de Pesquisa (RQs):
\begin{itemize}
    \item \textbf{RQ1 (Eficácia de Formalização):} \textit{O ciclo reflexivo guiado por JEV preserva ou aprimora a conformidade sintática ao template MASTeR e a fidelidade semântica em relação à geração em passo único (\textit{One-Shot}) e ao loop tradicional com LLM-as-a-Judge?}
    \item \textbf{RQ2 (Eficiência Operacional e Economia):} \textit{Qual é a redução obtida em latência de validação, volume de tokens gerados e custo financeiro por lote de 1.000 requisitos ao substituir a pilha de validadores generativos pelo JEV?}
    \item \textbf{RQ3 (Eficácia da Triagem por Complexidade):} \textit{O mecanismo de escalonamento por incerteza do JEV concentra os encaminhamentos humanos nos requisitos de maior complexidade estrutural (fator de divisão $\ge 2$), elevando a qualidade média do subconjunto aceito autonomamente?}
\end{itemize}

\subsection{Conjunto de Dados}
Utilizamos integralmente o benchmark público curado por Gro\ss er et al.~\cite{grosser2023comparative, grosser2024benchmarking} e disponibilizado no pacote suplementar de Ro\ss kothen et al.~\cite{rosskothen2026zenodo}. O conjunto compreende \textbf{249 requisitos fonte} em linguagem natural provenientes de cinco especificações heterogêneas: \textit{FLEX} (18 reqs., satélite), \textit{ECSS} (33 reqs., normas espaciais europeias), \textit{CSE} (25 reqs., certificação de motores aeronáuticos, com maior comprimento médio de $40{,}2 \pm 15{,}2$ palavras), \textit{TSS} (63 reqs., especificação técnica ferroviária) e \textit{EVS} (110 reqs., sistema eletrônico de votação, mais curtos e estruturados, com média de $20{,}2 \pm 8{,}3$ palavras). No total, especialistas humanos converteram esses 249 itens em \textbf{410 requisitos MASTeR} de referência (\textit{ground truth}), resultando em fator médio de divisão (\textit{split factor}) de $1{,}65$ ($168$ requisitos com fator $1$, $43$ com fator $2$ e $38$ com fator $3+$).

\subsection{Configurações Avaliadas}
Para assegurar tanto \textbf{validade interna} (isolamento estrito do mecanismo validador mantendo o gerador constante) quanto \textbf{validade externa} (generalização entre arquiteturas abertas e comerciais), avaliamos três modos de pipeline cruzados sobre duas famílias de modelos geradores (\texttt{Gemini 3.5 Flash} via API oficial Google e \texttt{Qwen3-235B-A22B-Instruct-2507} via API compatível, idêntico ao modelo primário de Ro\ss kothen et al.~\cite{rosskothen2026converting}), totalizando seis configurações experimentais controladas somadas às quatro linhas de base oficiais publicadas pelos autores:
\begin{enumerate}
    \item \textbf{\texttt{OneShot}:} Geração estruturada direta sem validadores ($0$ iterações extras), estabelecendo o limite inferior de custo e latência;
    \item \textbf{\texttt{LLM-Judge}:} Replicação fiel do loop reflexivo da Seção~III de Ro\ss kothen et al.~\cite{rosskothen2026converting} ($M = 3$ iterações máximas, limiar INCOSE de $0{,}7$), contabilizando individualmente tokens e latência do gerador, \textit{Template Validator}, \textit{INCOSE Validator} e \textit{Feedback Filter};
    \item \textbf{\texttt{JEV-Cascade}:} Pipeline proposto operando com \texttt{jev-latest} via SDK oficial \texttt{typesafe-sdk} ($M = 3$, $\tau_{\text{rej}} = 0{,}40$, $\tau_{\text{acc}} = 0{,}75$).
\end{enumerate}

\subsection{Métricas de Avaliação}
\subsubsection{Qualidade Sintática e Semântica}
Seguindo as descobertas da RQ2 de Ro\ss kothen et al.~\cite{rosskothen2026converting}, adotamos as métricas automáticas com maior correlação frente a especialistas humanos, todas normalizadas em $[0, 1]$: (a) \textbf{\texttt{MasterTemplateF1}} e \textbf{\texttt{MasterConditionF1}}, que calculam o $F_1$ de multiconjuntos entre os sub-templates MASTeR gerados e o \textit{ground truth}; (b) \textbf{\texttt{BLEUScore}} suavizado em nível de sentença; e, complementarmente quando habilitados modelos densos, (c) \textbf{\texttt{BERTScore}}~\cite{zhang2020bertscore} e \textbf{\texttt{EmbeddingMoversDistance} (EMD)}~\cite{kusner2015word}.

\subsubsection{Eficiência Operacional e Calibração}
Instrumentamos todas as chamadas assíncronas para registrar: (a) \textbf{latência média de validação} (em ms por requisito); (b) \textbf{consumo médio de tokens} de entrada e saída; (c) \textbf{custo estimado por 1.000 requisitos} (em USD); e (d) \textbf{taxas de roteamento} (\textit{Fast Accept Rate}, \textit{Escalation Rate} e o diferencial de qualidade $\Delta F_1$ entre o subconjunto aceito autonomamente e o subconjunto escalado para revisão humana). Para comparações pareadas entre configurações sobre os mesmos 249 requisitos, aplicamos o teste não paramétrico de postos sinalizados de Wilcoxon ($\alpha = 0{,}05$).

\section{Resultados Preliminares, Discussão e Ameaças à Validade}

\subsection{Diagnóstico Empírico das Linhas de Base (Ro\ss kothen et al., 2026)}
Antes da execução dos ciclos com JEV, reprocessamos os 996 registros oficiais ($249 \times 4$ configurações) disponibilizados por Ro\ss kothen et al.~\cite{rosskothen2026zenodo} para estabelecer a linha de base quantitativa do estudo. A Tabela~\ref{tab:baseline_paper} sintetiza o desempenho global e estratificado pelo fator de divisão humano (\textit{Split Factor} -- SF).

\begin{table}[t]
\centering
\caption{Desempenho das Configurações de Referência de Ro\ss kothen et al.~\cite{rosskothen2026zenodo} sobre os 249 Requisitos Fonte}
\label{tab:baseline_paper}
\resizebox{\columnwidth}{!}{%
\begin{tabular}{lcccccc}
\toprule
Configuração & \texttt{CorpBERT} $F_1$ & EMD & MASTeR $F_1$ & Cond. $F_1$ & BLEU & Composto \\
\midrule
\texttt{minimal-qwen235} (Sem Juiz) & \textbf{0,834} & \textbf{0,770} & 0,561 & \textbf{0,758} & \textbf{0,460} & \textbf{0,811} \\
\texttt{full-qwen235} (LLM-Judge)   & 0,823          & 0,762          & 0,553 & 0,742          & 0,423          & 0,797 \\
\texttt{minimal-qwen3-next}         & 0,822          & 0,758          & \textbf{0,592} & 0,725 & 0,412          & 0,801 \\
\texttt{full-gpt-5} (LLM-Judge)     & 0,787          & 0,725          & 0,569 & 0,656          & 0,273          & 0,753 \\
\midrule
\multicolumn{7}{l}{\textit{Degradação em \texttt{full-qwen235} conforme a complexidade estrutural (Split Factor):}} \\
\quad SF $= 1$ ($n=168$)            & 0,864          & 0,785          & 0,601 & ---            & ---            & --- \\
\quad SF $= 2$ ($n=43$)             & 0,766          & 0,734          & 0,478 & ---            & ---            & --- \\
\quad SF $\ge 3$ ($n=38$)           & 0,708          & 0,691          & 0,427 & ---            & ---            & --- \\
\bottomrule
\end{tabular}%
}
\end{table}

Os dados da Tabela~\ref{tab:baseline_paper} confirmam empiricamente as duas hipóteses motivadoras deste estudo:
\begin{enumerate}
    \item \textbf{Evidência de \textit{over-editing} no juiz generativo:} Ao comparar \texttt{minimal-qwen235} com \texttt{full-qwen235}, observa-se que a ativação do loop com juízes LLM generativos reduziu ligeiramente todas as métricas automáticas principais (ex.: escore composto caiu de $0{,}811$ para $0{,}797$), apesar de multiplicar o número de chamadas de inferência. Isso ocorre porque validadores não calibrados emitem críticas estilísticas sobre requisitos simples que já estavam corretos na primeira geração;
    \item \textbf{Colapso em requisitos multi-atômicos:} Em requisitos compostos com $\text{SF} \ge 3$, o \texttt{MasterTemplateF1} despenca de $0{,}601$ para $0{,}427$ ($-28{,}9\%$), demonstrando a necessidade crítica de um mecanismo calibrado de escalonamento humano (\textit{Path C} da nossa proposta).
\end{enumerate}

\subsection{Avaliação Comparativa do JEV-Cascade (RQ1--RQ3)}
A Tabela~\ref{tab:jev_results} consolida os resultados obtidos pelas execuções controladas do repositório experimental (\texttt{OneShot}, \texttt{LLM-Judge} e \texttt{JEV-Cascade}).
% TODO: Após rodar `scripts/run_experiment.py` e `scripts/evaluate_and_export_latex.py`, os números reais serão injetados automaticamente na tabela abaixo e discutidos nos parágrafos de RQ1, RQ2 e RQ3.

\begin{table}[t]
\centering
\caption{Comparação Experimental entre One-Shot, Loop LLM-as-a-Judge e JEV-Cascade ($N = 249$)}
\label{tab:jev_results}
\resizebox{\columnwidth}{!}{%
% Tabela inicial de espaço reservado (sobrescrita automaticamente por scripts/evaluate_and_export_latex.py)
\begin{tabular}{lcccccc}
\toprule
Configuração & MASTeR Main $F_1$ & Cond. $F_1$ & BLEU & Lat. Val. (ms) & Custo/1k (\$) & Escalon. (\%) \\
\midrule
\texttt{gemini\_oneshot}      & \textit{--} & \textit{--} & \textit{--} & $0{,}0$   & \textit{--} & $0{,}0\%$ \\
\texttt{gemini\_llm\_judge}   & \textit{--} & \textit{--} & \textit{--} & \textit{--} & \textit{--} & $0{,}0\%$ \\
\texttt{gemini\_jev\_cascade} & \textit{--} & \textit{--} & \textit{--} & \textit{--} & \textit{--} & \textit{--} \\
\texttt{qwen235\_jev\_cascade}& \textit{--} & \textit{--} & \textit{--} & \textit{--} & \textit{--} & \textit{--} \\
\bottomrule
\end{tabular}

}
\end{table}

\textbf{Análise esperada e guia de interpretação das RQs:}
\begin{itemize}
    \item \textbf{Sobre a RQ1 (Eficácia):} Avalia-se se o \texttt{JEV-Cascade} preserva os altos índices de \texttt{OneShot} nos itens aprovados via \textit{Fast Accept} enquanto recupera falhas gramaticais estritas (ausência de verbo modal ou prefixo condicional incorreto) por meio do mapeamento determinístico de reparo;
    \item \textbf{Sobre a RQ2 (Economia e Latência):} Diferentemente do modo \texttt{LLM-Judge}, no qual cada candidato aciona geração autorregressiva de comentários e uma etapa adicional de filtragem, cada auditoria JEV executa em um único passe paralelo sem tokens de saída, reduzindo substancialmente a latência acumulada de validação e o custo por milhar de requisitos;
    \item \textbf{Sobre a RQ3 (Calibração da Cascata):} Compara-se o \texttt{MasterTemplateF1} médio do subconjunto aceito autonomamente contra o subconjunto escalado ($\Delta F_1$), verificando se a zona de incerteza ($\tau_{\text{rej}} \le p_j < \tau_{\text{acc}}$) concentra efetivamente os casos com $\text{SF} \ge 2$.
\end{itemize}

\subsection{Ameaças à Validade}
Seguindo as diretrizes metodológicas de Runeson e H\"ost~\cite{runeson2009guidelines}, destacamos quatro categorias de ameaças à validade e suas respectivas mitigações:
\begin{itemize}
    \item \textbf{Validade de Construto:} As métricas automáticas comparam as formalizações geradas contra uma única reescrita humana de referência~\cite{grosser2023comparative}. Como demonstrado por Ro\ss kothen et al.~\cite{rosskothen2026converting}, existem múltiplas maneiras válidas de decompor requisitos complexos no padrão MASTeR. Mitigamos essa ameaça combinando métricas estruturais de multiconjunto (\texttt{MasterTemplateF1}) com métricas léxico-semânticas e analisando separadamente os estratos de complexidade;
    \item \textbf{Validade Interna:} Alterar simultaneamente o LLM gerador e o mecanismo validador impediria isolar a contribuição do JEV. Mitigamos esse risco preservando integralmente os prompts Jinja2 originais~\cite{rosskothen2026zenodo} e executando os três modos (\texttt{OneShot}, \texttt{LLM-Judge} e \texttt{JEV-Cascade}) dentro da mesma família de gerador (\texttt{Gemini 2.5 Flash}), além de suportar execução direta com \texttt{Qwen3-235B};
    \item \textbf{Validade Externa:} O estudo concentra-se no sistema de templates MASTeR sobre 249 requisitos de cinco projetos (predominantemente sistemas críticos e embarcados). Embora representativo de cenários industriais como PREEvision, a generalização para outros padrões semiformais (como EARS) ou requisitos ágeis (\textit{User Stories}) requer investigação adicional;
    \item \textbf{Confiabilidade:} Modelos em nuvem podem sofrer variações de latência de rede e atualizações de versão. Para garantir reprodutibilidade estrita, fixamos temperatura igual a zero ($T = 0{,}0$) com decodificação estruturada, persistimos telemetria individual por chamada em arquivos JSONL e disponibilizamos todo o pipeline de avaliação e testes automatizados.
\end{itemize}

\section{Conclusão e Trabalhos Futuros}

Este trabalho apresentou o \textit{JEV-Cascade Reflection Loop}, uma arquitetura híbrida que integra LLMs geradores a um modelo discriminativo de decisão rápida (\textit{System One} -- JEV) para converter requisitos em linguagem natural irrestrita em templates semiformais MASTeR. Ao substituir a sequência de validadores autorregressivos (\textit{LLM-as-a-Judge}) e o filtro auxiliar de feedback por uma única consulta paralela tipada (\texttt{Choice}, \texttt{Noul} e \texttt{Score}) condicionada simultaneamente ao requisito fonte e ao candidato, a abordagem ataca diretamente o custo operacional e o fenômeno de \textit{over-editing} observados nos fluxos reflexivos recentes da literatura.

Além de eliminar por construção a necessidade de um LLM dedicado à filtragem de feedback não acionável, o uso de probabilidades calibradas viabiliza uma política explícita de governança humano-no-ciclo (\textit{Fast Accept}, \textit{Guided Rewrite} e \textit{Escalate}), essencial para especificações industriais de grande porte onde requisitos compostos permanecem desafiadores para automação plena. Como trabalhos futuros, pretende-se (i) calibrar dinamicamente os limiares $(\tau_{\text{rej}}, \tau_{\text{acc}})$ via predição conformal sobre avaliações de especialistas, (ii) estender o catálogo de primitivas JEV para outros templates como EARS e (iii) investigar a correlação direta do escore ordinal do JEV como métrica substituta de avaliação humana em escala.


\bibliographystyle{IEEEtran}
\bibliography{references}

\end{document}
